\documentclass[11pt,a4paper]{article}
\usepackage[utf8]{utf8}
\usepackage[margin=0.8in]{geometry}
\usepackage{amsmath,amssymb}
\usepackage{enumitem}
\usepackage{booktabs}
\usepackage{titlesec}
\usepackage{xcolor}
\usepackage{listings}

% Setup clean styling for Python code blocks
\definecolor{codegray}{rgb}{0.5,0.5,0.5}
\definecolor{codepurple}{rgb}{0.58,0,0.82}
\definecolor{backcolour}{rgb}{0.96,0.96,0.96}

\lstdefinestyle{pystyle}{
	backgroundcolor=\color{backcolour},   
	commentstyle=\color{codegray},
	keywordstyle=\color{magenta},
	numberstyle=\tiny\color{codegray},
	stringstyle=\color{codepurple},
	basicstyle=\ttfamily\footnotesize,
	breakatwhitespace=false,         
	breaklines=true,                 
	captionpos=b,                    
	keepspaces=true,                 
	showspaces=false,                
	showstringspaces=false,
	showtabs=false,                  
	tabsize=2
}

\lstset{style=pystyle}

\titleformat{\section}{\large\bfseries\color{blue!70!black}}{\thesection.}{0.5em}{}[\titlerule]
\titleformat{\subsection}{\bfseries}{\thesubsection}{0.5em}{}
\setlist[enumerate]{leftmargin=*,itemsep=2pt,topsep=2pt}
\setlist[itemize]{leftmargin=*,itemsep=2pt,topsep=2pt}

\title{\textbf{Machine Learning Practicum: Linear \& Logistic Regression Worksheets}}
\author{\textbf{Department of Computer Science \& Data Science}}
\date{}

\begin{document}
	
	\maketitle
	
	\noindent\textbf{Instructions:} Solve the following 20 lab tasks using Python (\texttt{pandas}, \texttt{numpy}, \texttt{scikit-learn}, \texttt{matplotlib}). Ensure all input matrices $X$ are formatted as 2D arrays or DataFrames.
	
	\vspace{0.5em}
	
	% =============================================================
	\section{Linear Regression Lab Tasks}
	% =============================================================
	
	\subsection{Task 1: Single Feature House Price Prediction}
	\textbf{Scenario:} Predict house price ($y$, in \$100k) based on living area ($X$, in 1000 sq ft).
	\begin{lstlisting}[language=Python]
		import pandas as pd
		df1 = pd.DataFrame({
			'SqFt_1000': [1.2, 1.5, 1.8, 2.1, 2.4, 2.7, 3.0, 3.3],
			'Price_100k': [2.5, 2.9, 3.4, 3.8, 4.2, 4.8, 5.1, 5.5]
		})
	\end{lstlisting}
	\textbf{Tasks:}
	\begin{enumerate}
		\item Fit a \texttt{LinearRegression} model predicting \texttt{Price\_100k} from \texttt{SqFt\_1000}.
		\item Report the model intercept ($\beta_0$) and slope ($\beta_1$).
		\item Predict the price of a 2,200 sq ft house ($X = 2.2$).
	\end{enumerate}
	
	\subsection{Task 2: Multi-Feature Exam Performance}
	\textbf{Scenario:} Predict exam score ($y$, 0--100) using study hours ($X_1$) and sleep hours ($X_2$).
	\begin{lstlisting}[language=Python]
		df2 = pd.DataFrame({
			'Study_Hours': [2, 3, 4, 5, 6, 7, 8, 9],
			'Sleep_Hours': [8, 7, 6, 8, 7, 6, 8, 7],
			'Exam_Score':  [55, 62, 68, 77, 83, 88, 95, 98]
		})
	\end{lstlisting}
	\textbf{Tasks:}
	\begin{enumerate}
		\item Train a multiple linear regression model.
		\item Interpret the coefficient for \texttt{Study\_Hours} holding \texttt{Sleep\_Hours} constant.
		\item Predict the score for a student with 6 hours of study and 8 hours of sleep.
	\end{enumerate}
	
	\subsection{Task 3: Salary Estimation by Experience}
	\textbf{Scenario:} Predict annual software engineer salary ($y$, in \$1,000s) based on years of experience ($X$).
	\begin{lstlisting}[language=Python]
		df3 = pd.DataFrame({
			'Years_Exp': [1, 2, 3, 4, 5, 6, 7, 8],
			'Salary_1k': [50, 56, 63, 70, 78, 85, 94, 102]
		})
	\end{lstlisting}
	\textbf{Tasks:}
	\begin{enumerate}
		\item Calculate $R^2$ score and Mean Squared Error (MSE).
		\item Plot the raw data points and the fitted regression line using \texttt{matplotlib}.
	\end{enumerate}
	
	\subsection{Task 4: Vehicle Depreciation Analysis}
	\textbf{Scenario:} Predict used car price ($y$, in \$1,000s) from vehicle age ($X_1$, years) and odometer reading ($X_2$, in 10k miles).
	\begin{lstlisting}[language=Python]
		df4 = pd.DataFrame({
			'Age_Years': [1, 2, 3, 4, 5, 6, 7, 8],
			'Miles_10k': [1.0, 2.2, 3.5, 4.8, 5.5, 7.0, 8.2, 9.5],
			'Price_1k':  [24, 21, 18, 15, 13, 10, 8, 5]
		})
	\end{lstlisting}
	\textbf{Tasks:}
	\begin{enumerate}
		\item Fit a multiple linear regression model.
		\item Identify which feature has a stronger negative impact on price using standardized coefficients (\texttt{StandardScaler}).
	\end{enumerate}
	
	\subsection{Task 5: Utility Bill Modeling}
	\textbf{Scenario:} Estimate winter heating bill ($y$, USD) from square footage ($X_1$), average outdoor temperature ($X_2$, $^\circ$F), and number of occupants ($X_3$).
	\begin{lstlisting}[language=Python]
		df5 = pd.DataFrame({
			'SqFt':       [1200, 1500, 1800, 2000, 2200, 1400, 1600, 1900],
			'Avg_Temp_F': [35, 40, 45, 50, 55, 38, 42, 48],
			'Occupants':  [2, 3, 2, 4, 3, 1, 2, 4],
			'Bill_USD':   [150, 165, 170, 210, 220, 145, 160, 200]
		})
	\end{lstlisting}
	\textbf{Tasks:}
	\begin{enumerate}
		\item Compute the model's $R^2$ and Adjusted $R^2$.
		\item Explain why Adjusted $R^2$ is preferred when working with multiple features.
	\end{enumerate}
	
	\pagebreak
	
	\subsection{Task 6: Crop Yield Optimization}
	\textbf{Scenario:} Predict crop yield ($y$, tons/acre) based on nitrogen fertilizer applied ($X_1$, kg/acre) and rainfall ($X_2$, mm).
	\begin{lstlisting}[language=Python]
		df6 = pd.DataFrame({
			'Fertilizer_kg': [10, 20, 30, 40, 50, 60, 70, 80],
			'Rainfall_mm':   [100, 110, 120, 130, 140, 150, 160, 170],
			'Yield_Tons':    [2.1, 2.8, 3.4, 4.0, 4.5, 5.1, 5.6, 6.0]
		})
	\end{lstlisting}
	\textbf{Tasks:}
	\begin{enumerate}
		\item Fit linear regression and calculate residual values $e_i = y_i - \hat{y}_i$.
		\item Plot predicted values $\hat{y}$ vs. residuals $e$ to check for homoscedasticity.
	\end{enumerate}
	
	\subsection{Task 7: Tech Hardware Pricing}
	\textbf{Scenario:} Predict laptop price ($y$, USD) using RAM ($X_1$, GB) and SSD Storage ($X_2$, GB).
	\begin{lstlisting}[language=Python]
		df7 = pd.DataFrame({
			'RAM_GB':    [8, 8, 16, 16, 32, 32, 16, 64],
			'SSD_GB':    [256, 512, 512, 1024, 1024, 2048, 256, 2048],
			'Price_USD': [500, 620, 850, 1100, 1450, 1900, 780, 2400]
		})
	\end{lstlisting}
	\textbf{Tasks:}
	\begin{enumerate}
		\item Fit a linear regression model and write down the prediction equation.
		\item Estimate the price for a laptop with 32 GB RAM and 512 GB SSD.
	\end{enumerate}
	
	\subsection{Task 8: Digital Marketing ROI}
	\textbf{Scenario:} Predict monthly sales revenue ($y$, in \$1,000s) from ad spend ($X_1$, in \$1,000s) and social media posts ($X_2$).
	\begin{lstlisting}[language=Python]
		df8 = pd.DataFrame({
			'Ad_Spend_1k':  [1.0, 1.5, 2.0, 2.5, 3.0, 3.5, 4.0, 4.5],
			'Social_Posts': [10, 15, 20, 22, 28, 30, 35, 40],
			'Revenue_1k':   [12, 16, 21, 25, 30, 34, 40, 44]
		})
	\end{lstlisting}
	\textbf{Tasks:}
	\begin{enumerate}
		\item Train the model and report Mean Squared Error (MSE).
		\item Calculate the feature correlation matrix using \texttt{df8.corr()}.
	\end{enumerate}
	
	\subsection{Task 9: Delivery Logistics}
	\textbf{Scenario:} Predict delivery time ($y$, minutes) based on route distance ($X_1$, miles) and drop-off stops ($X_2$).
	\begin{lstlisting}[language=Python]
		df9 = pd.DataFrame({
			'Distance_Miles': [2, 4, 6, 8, 10, 12, 14, 16],
			'Stops':          [1, 2, 2, 3, 4, 4, 5, 6],
			'Time_Min':       [15, 25, 32, 42, 53, 60, 71, 82]
		})
	\end{lstlisting}
	\textbf{Tasks:}
	\begin{enumerate}
		\item Fit a multiple linear regression model.
		\item Determine the expected delay per additional delivery stop holding distance constant.
	\end{enumerate}
	
	\subsection{Task 10: Restaurant Revenue Metrics}
	\textbf{Scenario:} Predict daily revenue ($y$, in \$1,000s) based on seating capacity ($X_1$) and daily foot traffic ($X_2$).
	\begin{lstlisting}[language=Python]
		df10 = pd.DataFrame({
			'Seats':        [30, 40, 50, 60, 70, 80, 90, 100],
			'Foot_Traffic': [150, 200, 280, 320, 400, 450, 520, 600],
			'Revenue_1k':   [2.1, 2.9, 3.8, 4.5, 5.4, 6.2, 7.1, 8.0]
		})
	\end{lstlisting}
	\textbf{Tasks:}
	\begin{enumerate}
		\item Fit \texttt{LinearRegression} and evaluate using $R^2$.
		\item Compare actual vs. predicted revenue for all 8 observations.
	\end{enumerate}
	
	\pagebreak
	
	% =============================================================
	\section{Logistic Regression Lab Tasks}
	% =============================================================
	
	\subsection{Task 1: Single Feature Pass/Fail Prediction}
	\textbf{Scenario:} Predict whether a student passes an exam ($y \in \{0, 1\}$) based on study hours ($X$).
	\begin{lstlisting}[language=Python]
		clf1 = pd.DataFrame({
			'Study_Hours': [0.5, 1.0, 1.5, 2.0, 2.5, 3.0, 3.5, 4.0, 4.5, 5.0],
			'Passed':      [0,   0,   0,   0,   0,   1,   0,   1,   1,   1]
		})
	\end{lstlisting}
	\textbf{Tasks:}
	\begin{enumerate}
		\item Fit a \texttt{LogisticRegression} model predicting \texttt{Passed}.
		\item Calculate the decision boundary where $P(\text{Passed}) = 0.5$.
		\item Predict pass probability for a student who studied $2.8$ hours.
	\end{enumerate}
	
	\subsection{Task 2: Credit Risk Assessment}
	\textbf{Scenario:} Predict loan approval ($1 = \text{Approved}, 0 = \text{Denied}$) based on credit score ($X_1$) and annual income ($X_2$, in \$1,000s).
	\begin{lstlisting}[language=Python]
		clf2 = pd.DataFrame({
			'Credit_Score': [580, 610, 640, 660, 700, 720, 750, 800],
			'Income_1k':    [35,  40,  45,  60,  55,  80,  95,  110],
			'Approved':     [0,   0,   0,   0,   1,   1,   1,   1]
		})
	\end{lstlisting}
	\textbf{Tasks:}
	\begin{enumerate}
		\item Train a logistic regression classifier.
		\item Output predicted probabilities for all samples using \texttt{model.predict\_proba(X)}.
		\item Calculate training set accuracy score.
	\end{enumerate}
	
	\subsection{Task 3: Email Spam Detection}
	\textbf{Scenario:} Classify an email as spam ($1$) or not spam ($0$) based on word count ($X_1$) and link count ($X_2$).
	\begin{lstlisting}[language=Python]
		clf3 = pd.DataFrame({
			'Word_Count': [50, 120, 450, 800, 30, 200, 600, 90],
			'Link_Count': [0,  1,   8,   12,  0,  2,   10,  5],
			'Is_Spam':    [0,  0,   1,   1,   0,  0,   1,   1]
		})
	\end{lstlisting}
	\textbf{Tasks:}
	\begin{enumerate}
		\item Fit logistic regression and display model intercept and coefficients.
		\item Generate a confusion matrix using \texttt{sklearn.metrics.confusion\_matrix}.
	\end{enumerate}
	
	\subsection{Task 4: Heart Disease Screening}
	\textbf{Scenario:} Predict high heart disease risk ($1$) vs. low risk ($0$) using age ($X_1$) and cholesterol levels ($X_2$, mg/dL).
	\begin{lstlisting}[language=Python]
		clf4 = pd.DataFrame({
			'Age':         [30, 40, 45, 50, 55, 60, 65, 70],
			'Cholesterol': [180, 195, 210, 240, 220, 260, 280, 300],
			'High_Risk':   [0,   0,   0,   1,   0,   1,   1,   1]
		})
	\end{lstlisting}
	\textbf{Tasks:}
	\begin{enumerate}
		\item Fit the model and compute Precision, Recall, and F1-score.
		\item Identify any false positive or false negative predictions.
	\end{enumerate}
	
	\subsection{Task 5: Customer Churn Prediction}
	\textbf{Scenario:} Predict customer churn ($1 = \text{Yes}, 0 = \text{No}$) using monthly charge ($X_1$, USD) and tenure ($X_2$, months).
	\begin{lstlisting}[language=Python]
		clf5 = pd.DataFrame({
			'Monthly_Charge': [25, 35, 65, 85, 45, 95, 105, 30],
			'Tenure_Months':  [24, 36, 6,  2,  18, 4,  1,   48],
			'Churned':        [0,  0,  1,  1,  0,  1,  1,   0]
		})
	\end{lstlisting}
	\textbf{Tasks:}
	\begin{enumerate}
		\item Train \texttt{LogisticRegression}.
		\item Convert log-odds coefficients to Odds Ratios using $\text{OR} = e^\beta$.
		\item Interpret how tenure affects churn probability.
	\end{enumerate}
	
	\pagebreak
	
	\subsection{Task 6: University Admission Model}
	\textbf{Scenario:} Predict admission ($1 = \text{Admitted}, 0 = \text{Rejected}$) using GPA ($X_1$, 0--4.0 scale) and SAT score ($X_2$, 400--1600).
	\begin{lstlisting}[language=Python]
		clf6 = pd.DataFrame({
			'GPA':      [2.8, 3.1, 3.4, 3.5, 3.7, 3.8, 3.9, 4.0],
			'SAT':      [1050, 1120, 1200, 1280, 1340, 1420, 1490, 1550],
			'Admitted': [0,    0,    0,    1,    1,    1,    1,    1]
		})
	\end{lstlisting}
	\textbf{Tasks:}
	\begin{enumerate}
		\item Fit model and plot predicted probability curve against GPA.
		\item Predict admission status for a student with GPA = $3.3$ and SAT = $1220$.
	\end{enumerate}
	
	\subsection{Task 7: Employee Promotion Likelihood}
	\textbf{Scenario:} Predict promotion status ($1 = \text{Promoted}, 0 = \text{Not Promoted}$) based on performance rating ($X_1$, 1--5) and tenure ($X_2$, years).
	\begin{lstlisting}[language=Python]
		clf7 = pd.DataFrame({
			'Rating':   [2, 3, 3, 4, 4, 5, 5, 5],
			'Years':    [1, 2, 5, 3, 6, 2, 4, 7],
			'Promoted': [0, 0, 0, 0, 1, 1, 1, 1]
		})
	\end{lstlisting}
	\textbf{Tasks:}
	\begin{enumerate}
		\item Fit \texttt{LogisticRegression} and print \texttt{classification\_report}.
		\item Evaluate overall prediction accuracy on the dataset.
	\end{enumerate}
	
	\subsection{Task 8: E-Commerce Purchase Intent}
	\textbf{Scenario:} Predict if a visitor buys an item ($1$) or leaves ($0$) using time on site ($X_1$, seconds) and pages viewed ($X_2$).
	\begin{lstlisting}[language=Python]
		clf8 = pd.DataFrame({
			'Time_Sec':  [20, 45, 90, 150, 210, 300, 420, 600],
			'Pages':     [1,  2,  3,  4,   5,   7,   8,   12],
			'Purchased': [0,  0,  0,  0,   1,   1,   1,   1]
		})
	\end{lstlisting}
	\textbf{Tasks:}
	\begin{enumerate}
		\item Standardize features using \texttt{StandardScaler} prior to fitting.
		\item Compute ROC-AUC score using \texttt{roc\_auc\_score}.
	\end{enumerate}
	
	\subsection{Task 9: Real Estate Deal Closure}
	\textbf{Scenario:} Classify if an offer is accepted ($1$) or rejected ($0$) based on discount off list price ($X_1$, \%) and days on market ($X_2$).
	\begin{lstlisting}[language=Python]
		clf9 = pd.DataFrame({
			'Discount_Pct': [2,  4,  6,  8,  10, 12, 15, 18],
			'Days_Market':  [12, 20, 45, 60, 30, 90, 80, 110],
			'Accepted':     [1,  1,  1,  0,  1,  0,  0,  0]
		})
	\end{lstlisting}
	\textbf{Tasks:}
	\begin{enumerate}
		\item Fit a logistic regression model.
		\item Identify which feature negatively impacts acceptance probability more strongly.
	\end{enumerate}
	
	\subsection{Task 10: Botanical Plant Survival}
	\textbf{Scenario:} Predict plant survival ($1 = \text{Survived}, 0 = \text{Died}$) based on daily sunlight ($X_1$, hours) and daily water ($X_2$, mL).
	\begin{lstlisting}[language=Python]
		clf10 = pd.DataFrame({
			'Sunlight_Hrs': [2, 3, 4, 5, 6, 7, 8, 9],
			'Water_mL':     [50, 80, 100, 150, 200, 250, 300, 350],
			'Survived':     [0,  0,  0,   1,   1,   1,   1,   1]
		})
	\end{lstlisting}
	\textbf{Tasks:}
	\begin{enumerate}
		\item Fit \texttt{LogisticRegression}.
		\item Predict survival class and probability for a plant receiving $4.5$ hours of sunlight and $120$ mL water.
	\end{enumerate}
	
\end{document}